\diarytitle{0040}{非同次オイラー方程式（コーシー・オイラー方程式）の初期値問題（重根・非同次項 $\ln x$）}

オイラー方程式（コーシー・オイラー方程式）$x^{2}y'' + axy' + by = f(x)$ は $x = e^{t}$（$x>0$）と置換すると定数係数線形微分方程式に帰着できる。$t=\ln x$ とすると $x\dfrac{dy}{dx}=\dfrac{dy}{dt}$、$x^{2}\dfrac{d^{2}y}{dx^{2}}=\dfrac{d^{2}y}{dt^{2}}-\dfrac{dy}{dt}$ であるから、$D=\dfrac{d}{dt}$ と略記すると方程式は $\dfrac{d^{2}y}{dt^{2}}+(a-1)\dfrac{dy}{dt}+by=f(e^{t})$ という $t$ に関する定数係数線形微分方程式になる。

$t = \ln x$ とおくと、上記の変換により
\[
\frac{d^{2}y}{dt^{2}} - 2\frac{dy}{dt} + y = t
\]
となる（$a=-1,\ b=1$ より $a-1=-2$、右辺は $f(e^{t}) = \ln(e^{t}) = t$）。

まず斉次方程式の特性方程式は $m^{2}-2m+1=0$ より $(m-1)^{2}=0$、重根 $m=1$。よって
\[
y_{h} = (C_{1}+C_{2}t)e^{t}
\]
次に特解を $y_{p} = At+B$（$t=0,1$ は特性方程式の根でないので多項式のまま仮定できる）とおくと、$y_{p}'=A,\ y_{p}''=0$ であるから
\[
0 - 2A + (At+B) = t
\quad\Longrightarrow\quad
At + (B-2A) = t
\]
係数比較により $A=1,\ B=2A=2$。よって $y_{p} = t+2$。

したがって $t$ に関する一般解は $y = (C_{1}+C_{2}t)e^{t} + t + 2$。$t=\ln x,\ e^{t}=x$ を代入して $x$ に戻すと
\[
y = C_{1}x + C_{2}x\ln x + \ln x + 2
\]
初期条件を適用する。$y' = C_{1} + C_{2}(\ln x + 1) + \dfrac{1}{x}$ である。

$x=1$（$\ln 1 = 0$）で
\[
y(1) = C_{1} + 2 = 1 \implies C_{1} = -1
\]
\[
y'(1) = C_{1} + C_{2} + 1 = 0 \implies -1 + C_{2} + 1 = 0 \implies C_{2} = 0
\]
よって
\[
\boxed{\; y = 2 - x + \ln x \;}
\]
（検算: $y'=-1+\dfrac{1}{x}$, $y''=-\dfrac{1}{x^{2}}$ より
$x^{2}y'' - xy' + y = -1 - x\left(-1+\dfrac{1}{x}\right) + (2-x+\ln x) = -1 + x - 1 + 2 - x + \ln x = \ln x$ で方程式を満たす。
また $y(1)=2-1+0=1$、$y'(1)=-1+1=0$ で初期条件も満たす。）
